\clearpage
\begin{center}
{\Large BESTIMMUNG}\par\vspace{8mm}
{\small DES}\par\vspace{8mm}
{\LARGE\textls[100]{BREITENUNTERSCHIEDS}}\par\vspace{8mm}
{\small ZWISCHEN DEN}\par\vspace{7mm}
{\large STERNWARTEN VON GÖTTINGEN UND ALTONA}\par\vspace{7mm}
{\small DURCH}\par\vspace{7mm}
{\normalsize BEOBACHTUNGEN MIT DEM RAMSDENSCHEN ZENITSEKTOR}\par\vspace{7mm}
{\small VON}\par\vspace{7mm}
{\large\textls[100]{CARL FRIEDRICH GAUSS,}}\par\smallskip
{\scriptsize
RITTER DES GUELPHEN- UND DANNEBROG-ORDENS; KÖNIGLICH GROSSBRITANNISCH-HANNOVERSCHER HOFRAT;\\
PROFESSOR DER ASTRONOMIE UND DIREKTOR DER STERNWARTE IN GÖTTINGEN;\\
MITGLIED DER AKADEMIEN UND SOZIETÄTEN VON BERLIN, KOPENHAGEN, EDINBURGH, GÖTTINGEN,
LONDON, MÜNCHEN, NEAPEL, PARIS, PETERSBURG UND STOCKHOLM SOWIE DER
AMERIKANISCHEN, ITALIENISCHEN, KURLÄNDISCHEN, LONDONER ASTRONOMISCHEN U. A.
}\par\vspace{15mm}
\rule{.35\textwidth}{.4pt}\par\medskip
{\large\textls[170]{GÖTTINGEN,}}\par\medskip
BEI VANDENHOECK UND RUPRECHT.\par\medskip
1828.
\end{center}
\vfill\noindent IX.\hfill 1\par

\clearpage
\null

\clearpage
\begin{center}
{\LARGE\textls[170]{INHALT.}}\par\medskip
\rule{.09\textwidth}{.4pt}
\end{center}
\noindent Einleitung\dotfill S.\quad 5\par
\noindent\makebox[8mm][r]{I.}\quad Die beobachteten Sterne.\dotfill \quotedblbase\quad 8\par
\noindent\makebox[8mm][r]{II.}\quad Die Beobachtungen\dotfill \quotedblbase\quad 10\par
\noindent\makebox[8mm][r]{III.}\quad Ergebnisse\dotfill \quotedblbase\quad 29\par
\noindent\hspace*{12mm}\begin{minipage}{\dimexpr\textwidth-24mm\relax}
Einfachste Kombination der Beobachtungen zur Bestimmung des Breitenunterschieds, Art.\ 1.
Genauigkeit der Beobachtungen, Art.\ 2. Kollimationsfehler, Art.\ 3.
Absolut vorteilhafteste Kombination der Beobachtungen, Art.\ 4---7.
Berücksichtigung der unregelmäßigen Teilungsfehler, Art.\ 8---11. Lage der Beobachtungsplätze, Art.\ 12.
Bestimmung der absoluten Polhöhe der Göttinger Sternwarte aus Beobachtungen des Polarsterns
am \textsc{Reichenbach}schen Meridiankreis, Art.\ 13---17. Endergebnis der hannoverschen
Gradmessung, Art.\ 18---20. Vergleich der Deklinationen der beobachteten Zenitsterne mit den
Bestimmungen von \textsc{Bradley} und \textsc{Piazzi}, Art.\ 21.
\end{minipage}\par
\noindent\makebox[8mm][r]{IV.}\quad Breitenbestimmung der Sternwarte Seeberg\dotfill \quotedblbase\quad 52\par
\noindent Zusatz zu [Art.\ 20] S.\ 48\dotfill \quotedblbase\quad 56\par
\begin{center}\rule{.23\textwidth}{.4pt}\end{center}
\vfill\begin{flushright}$1^{\ast}$\end{flushright}

\clearpage
\null

\clearpage
\vspace*{12mm}
\begin{center}
{\large BESTIMMUNG DES BREITENUNTERSCHIEDS}\par\medskip
{\small ZWISCHEN}\par\medskip
{\large DEN STERNWARTEN VON GÖTTINGEN UND ALTONA\\
DURCH BEOBACHTUNGEN\\
MIT DEM RAMSDENSCHEN ZENITSEKTOR.}\par\medskip
\rule{.12\textwidth}{.4pt}\par\bigskip
{\large EINLEITUNG.}
\end{center}

Durch die Dreieckskette, die ich in den Jahren 1821---1824 durch das Königreich Hannover
entlang des Göttinger Meridians geführt habe, sind die Sternwarten von Göttingen und Altona
mit größter Genauigkeit trigonometrisch miteinander verbunden. Diese Messungen werden
später ausführlich veröffentlicht. Hier sei nur bemerkt, dass die absoluten Größen
auf der von Herrn Professor \textsc{Schumacher} in Holstein mit äußerster Genauigkeit
gemessenen Basis beruhen, mit der das Dreieckssystem durch die Seite Hamburg-Hohenhorn
verbunden ist. Die Orientierung stützt sich auf die Beobachtungen mit dem Göttinger
Passagenfernrohr, da die Sternwarte und das nördliche Meridianzeichen selbst Dreieckspunkte
sind. Die Sternwarten von Göttingen und Altona liegen durch einen merkwürdigen Zufall
bis auf weniger als eine Hausbreite auf demselben Meridian. Obwohl die absoluten Polhöhen
durch Beobachtungen mit fest aufgestellten Meridianinstrumenten bestimmt sind, war es
doch wichtig, den Breitenunterschied noch auf eine andere Weise mit ein und demselben
Instrument zu bestimmen, und ich hatte das Glück,

\clearpage
\noindent\makebox[10mm][l]{6}\parbox[t]{\dimexpr\textwidth-10mm\relax}{\centering\scriptsize BESTIMMUNG DES\,\quotedblbase\,BREITENUNTERSCHIEDS ZWISCHEN GÖTTINGEN UND ALTONA.}\par
\noindent dazu den ausgezeichneten Zenitsektor von \textsc{Ramsden} verwenden zu können,
der bekanntlich bei der englischen Gradmessung einem ähnlichen Zweck gedient hat.
Die Beobachtungen, die ich damit im Frühjahr 1827 angestellt habe, und ihre Ergebnisse
sind der Hauptgegenstand dieser Schrift.

Wenn viele Sterne nacheinander zu beobachten sind, lassen sich die Beobachtungen mit diesem
Instrument kaum ohne die Hilfe eines geübten Assistenten ausführen. Herr Professor
\textsc{Schumacher} hatte daher die Güte, Herrn Ingenieurleutnant von
\textsc{Nehus}[\textsuperscript{*})] mit Genehmigung Seiner Majestät des Königs von Dänemark
für die Beobachtungen an beiden Orten damit zu beauftragen. Dieser sehr geschickte Beobachter
übernahm fortlaufend das Ablesen der Mikrometerschraube und das Einstellen des Lotfadens.
Ich selbst beobachtete die Durchgänge an den Meridianfäden und stellte den zum Meridian
senkrechten Faden auf die Sterne ein. Nur in den beiden ersten Beobachtungsnächten in Altona
übernahm ein anderer Assistent jene Aufgabe. Deshalb habe ich diese Beobachtungen nicht
aufgenommen, zumal die Erfahrung bestätigte, dass verschiedene Personen die Halbierung
der Punkte durch den Lotfaden unterschiedlich einschätzten.

Das Instrument ist durch die ausführliche Beschreibung von \textsc{Mudge} hinreichend
bekannt. In Göttingen konnte es in der Sternwarte selbst unter dem östlichen Meridianspalt
aufgestellt werden. In Altona war das nicht möglich. Daher wurde es im Garten von Herrn
Professor \textsc{Schumacher}, in dem die dortige Sternwarte selbst steht, unter demselben
Beobachtungszelt aufgestellt, das \textsc{Mudge} in England benutzt hatte. Die Stabilität der
Aufstellung auf eingerammten Pfählen ließ nichts zu wünschen übrig. Die Nivellierung der
Vertikalachse wurde täglich überprüft, und gewöhnlich war fast nichts zu ändern. Dasselbe
gilt für die Horizontalachse.

Um die Ebene des Teilkreises in den Meridian zu bringen, wurde in Göttingen das südliche
Meridianzeichen benutzt. Es steht zwar im Meridian des westlichen Spalts, doch ließ sich sein
Azimut am Standort des Sektors mit größter Genauigkeit berechnen. In Altona ließ sich kein
ähnliches Verfahren anwenden. Dort wurde der Teilkreis zunächst mithilfe der bekannten
absoluten Zeit anhand eines kulminierenden Sterns sehr nahe in den

\vfill\noindent\rule{.24\textwidth}{.4pt}\par
{\footnotesize\noindent [\textsuperscript{*}) Handschriftliche Bemerkung:] \textsc{Nehus} starb am 17. April 1844 an zurückgetretener Grippe.\par}

\clearpage
\noindent\mbox{}\hfill EINLEITUNG.\hfill 7\par
\noindent Meridian gebracht. Durch Beobachtung mehrerer Sterne über die ganze Ausdehnung
des Teilkreises ließ sich dann leicht die noch nötige kleine Korrektur der Einstellung
bestimmen. Wie bereits erwähnt, wurden jede Nacht auch die Durchgänge der am Sektor
kulminierenden Sterne an den Meridianfäden beobachtet, und die Rektaszensionen der Sterne
waren bekannt. Damit ließ sich die richtige Ausrichtung im Meridian fortwährend zuverlässig
überprüfen. Nur einmal war eine geringfügige Nachstellung nötig. In der Regel wechselte
von einer Nacht zur nächsten die Stellung des Teilkreises zwischen Ost und West. Nur gegen
Ende der Beobachtungen wurde von dieser Regel gelegentlich abgewichen, um die Zahl der
Beobachtungen auf beide Lagen annähernd gleich zu verteilen. Dann wurde der Sektor in
einer Nacht einmal oder mehrmals gewendet.

Die Anzeigen des Barometers sowie des Innen- und Außenthermometers wurden jede Nacht
mindestens dreimal aufgezeichnet, zu Beginn, in der Mitte und am Ende der Beobachtungen.
Nach dem Vorbild von \textsc{Mudge} wurde ebenso der Temperaturunterschied zwischen dem
oberen und unteren Ende des Sektors aufgezeichnet, da er den Teilkreis und den Radius
in ungleichem Verhältnis verändert. Dass im Übrigen jede weitere durch die Bauart des
Instruments vorgeschriebene Vorsichtsmaßnahme sorgfältig beachtet wurde, braucht wohl
nicht eigens erwähnt zu werden. Dazu gehörte etwa, das Wassergefäß, in das das Lot hängt,
ausreichend gefüllt zu halten und möglichst dieselben Gewindegänge der Mikrometerschraube
zu benutzen. Bei jedem Stern wurde der Lotfaden unabhängig von der vorhergehenden Einstellung
neu auf den oberen Punkt (den Mittelpunkt des Gradbogens) eingestellt. Die Einstellung
auf den nächsten Teilungspunkt (oder die beiden nächsten) wurde in der Regel mehrmals
wiederholt. Aus den verschiedenen Ablesungen der Mikrometerschraube, die meistens bis auf
wenige Dezimalbruchteile einer Sekunde übereinstimmten, wurde der Mittelwert gebildet.

\begin{center}\rule{.24\textwidth}{.4pt}\end{center}

\clearpage
\vspace*{18mm}
\begin{center}
{\large I.}\par\medskip
{\Large DIE BEOBACHTETEN STERNE.}
\end{center}

Anfangs hatte ich 38 Sterne in für die Beobachtung geeigneten Lagen ausgewählt.
Gegen Ende der Beobachtungen in Göttingen fügte ich noch fünf weitere hinzu, weil ich
befürchtete, ungünstiges Wetter könnte den Abschluss der Beobachtungen in Altona so weit
verzögern, dass ein beträchtlicher Teil der zuerst ausgewählten Sterne wegen ihrer dann
tagsüber eintretenden Kulmination nicht oft genug beobachtet werden könnte. Diese Befürchtung
bestätigte sich jedoch nur in geringem Maße, und nur ein einziger Stern wurde in Altona bloß
einseitig beobachtet. Ich gebe hier die mittlere Stellung dieser Sterne, reduziert auf den
Anfang des Jahres 1827. Die Deklinationen sind die Ergebnisse der Beobachtungen mit dem
Zenitsektor selbst. Die Rektaszensionen, deren genaueste Bestimmung für den vorliegenden
Zweck unwesentlich ist, beruhen meist nur auf einer einmaligen Beobachtung am Meridiankreis,
deren Reduktion Herr \textsc{von Heiligenstein} freundlicherweise berechnet hat.
Der Bequemlichkeit halber bezeichne ich die Sterne mit fortlaufenden Nummern. Nr.\ 8, 13, 15
und 31 sind Doppelsterne. Beim ersten wurde stets der nachfolgende Stern eingestellt, bei
den beiden folgenden die Mitte. Bei Nr.\ 31 ist der Begleitstern so lichtschwach, dass er
im Fernrohr des Sektors selbst bei versuchsweise verdunkeltem Gesichtsfeld stets unsichtbar
blieb. Herr Professor \textsc{Schumacher} bemerkte ihn dagegen im lichtstärkeren Fernrohr
des \textsc{Reichenbach}schen Meridiankreises sofort, obwohl uns damals nicht bekannt war,
dass andere Astronomen diesen Stern bereits als Doppelstern erkannt hatten.

\clearpage
\noindent\parbox[t]{\dimexpr\textwidth-10mm\relax}{\centering\scriptsize BESTIMMUNG DES BREITENUNTERSCHIEDS USW.\quad I.\quad DIE BEOBACHTETEN STERNE.}\makebox[10mm][r]{9}\par
\begin{center}
\begin{minipage}[t]{.49\textwidth}\centering
{\scriptsize\setlength{\tabcolsep}{1.2pt}\renewcommand{\arraystretch}{1.2}\begin{tabular}{r|l|rrr|rrr|}
\multicolumn{2}{c|}{Bezeichnung}&\multicolumn{3}{c|}{Rektasz. 1827}&\multicolumn{3}{c|}{Dekl. 1827}\\\hline
1&24 Canum&$13^{\mathrm h}$&$27^{\mathrm m}$&$22^{\mathrm s}{,}39$&$49^\circ$&$54'$&$11''{,}62$\\
2&83 Ursae&&$34$&$9{,}77$&$55$&$33$&$34{,}98$\\
3&$\eta$ Ursae&&$40$&$42{,}86$&$50$&$10$&$46{,}20$\\
4&86 Ursae&&$47$&$28{,}48$&$54$&$34$&$55{,}55$\\
5&&&$50$&$45{,}55$&$55$&$25$&$59{,}22$\\
6&P. 13. 289&&$55$&$19{,}44$&$46$&$35$&$39{,}14$\\
7&13 Bootis&$14$&$1$&$49{,}02$&$50$&$16$&$43{,}44$\\
8&$\varkappa$ Bootis seq.&&$7$&$16{,}93$&$52$&$36$&$7{,}47$\\
9&P. 14. 56&&$12$&$5{,}39$&$56$&$13$&$36{,}51$\\
10&$\theta$ Bootis&&$19$&$18{,}40$&$52$&$39$&$12{,}05$\\
11&P. 14. 131&&$27$&$50{,}42$&$53$&$39$&$33{,}02$\\
12&P. 14. 164&&$35$&$24{,}18$&$52$&$58$&$55{,}66$\\
13&39 Bootis med.&&$43$&$48{,}37$&$49$&$26$&$9{,}46$\\
14&P. 14. 235&&$50$&$38{,}71$&$50$&$20$&$22{,}49$\\
15&44 Bootis med.&&$58$&$5{,}30$&$48$&$19$&$52{,}47$\\
16&&$15$&$7$&$8{,}12$&$49$&$13$&$46{,}68$\\
17&P. 15. 39&&$10$&$33{,}95$&$51$&$34$&$53{,}72$\\
18&&&$15$&$0{,}55$&$52$&$35$&$5{,}87$\\
19&&&$21$&$50{,}88$&$54$&$37$&$34{,}87$\\
20&&&$30$&$42{,}13$&$54$&$29$&$54{,}30$\\
21&&&$38$&$9{,}07$&$52$&$54$&$37{,}09$\\
22&&&$42$&$14{,}32$&$46$&$16$&$4{,}70$\\
\end{tabular}}\end{minipage}\hfill
\begin{minipage}[t]{.49\textwidth}\centering
{\scriptsize\setlength{\tabcolsep}{1.2pt}\renewcommand{\arraystretch}{1.2}\begin{tabular}{r|l|rrr|rrr}
\multicolumn{2}{c|}{Bezeichnung}&\multicolumn{3}{c|}{Rektasz. 1827}&\multicolumn{3}{c|}{Dekl. 1827}\\\hline
23&&$15^{\mathrm h}$&$48^{\mathrm m}$&$16^{\mathrm s}{,}24$&$56^\circ$&$20'$&$27''{,}00$\\
24&&&$54$&$10{,}23$&$50$&$22$&$39{,}41$\\
25&$\theta$ Draconis&&$58$&$39{,}70$&$59$&$1$&$47{,}14$\\
26&&$16$&$4$&$24{,}06$&$50$&$38$&$13{,}12$\\
27&P. 16. 33&&$7$&$24{,}92$&$46$&$20$&$17{,}04$\\
28&P. 16. 56&&$11$&$34{,}31$&$53$&$40$&$13{,}41$\\
29&&&$16$&$2{,}19$&$52$&$27$&$11{,}06$\\
30&&&$20$&$38{,}75$&$55$&$36$&$4{,}73$\\
31&&&$26$&$34{,}90$&$45$&$58$&$5{,}85$\\
32&16 Draconis&&$32$&$6{,}42$&$53$&$15$&$3{,}64$\\
33&&&$37$&$53{,}85$&$50$&$16$&$13{,}67$\\
34&&&$42$&$1{,}51$&$57$&$5$&$37{,}79$\\
35&&&$45$&$2{,}26$&$46$&$56$&$48{,}53$\\
36&P. 16. 253&&$49$&$21{,}51$&$46$&$49$&$22{,}01$\\
37&P. 16. 291&&$56$&$11{,}73$&$56$&$56$&$42{,}35$\\
38&P. 16. 310&$17$&$0$&$15{,}10$&$49$&$2$&$47{,}09$\\
39&P. 17. 20&&$4$&$23{,}12$&$58$&$29$&$49{,}33$\\
40&P. 17. 38&&$8$&$1{,}46$&$56$&$52$&$21{,}74$\\
41&74 Herculis&&$15$&$28{,}30$&$46$&$24$&$52{,}90$\\
42&P. 17. 120&&$20$&$53{,}56$&$57$&$10$&$13{,}11$\\
43&$\beta$ Draconis&&$26$&$31{,}90$&$52$&$25$&$57{,}91$\\
&&&&&&&\\
\end{tabular}}\end{minipage}
\end{center}
\begin{center}\rule{.24\textwidth}{.4pt}\end{center}

\clearpage
\vspace*{18mm}
\begin{center}
{\large II.}\par\medskip
{\Large DIE BEOBACHTUNGEN.}
\end{center}

Ein vollständiger Abdruck des Tagebuchs in seiner ursprünglichen Form hätte den Umfang
dieser Schrift mehr als verdoppelt und schien mir überflüssig. Ich gebe daher die
Beobachtungen sogleich nach Sternen geordnet wieder.

Die erste Spalte enthält die Zenitdistanzen, wie sie das Instrument geliefert hat,
das heißt die bloße Reduktion der Ablesung. Nördliche Zenitdistanzen gelten als positiv,
südliche als negativ.

Die zweite Spalte gibt die gemeinsame Wirkung der Refraktion und der ungleichen Ausdehnung
des Instruments infolge des Temperaturunterschieds zwischen dem oberen und dem unteren Teil
an. Die größten vorgekommenen Unterschiede waren $+1^\circ{,}2$ Réaum.\ (das obere
Thermometer zeigte mehr an) und $-0^\circ{,}6$. Um ein unverfälschtes Ergebnis zur
Beurteilung der gegenseitigen Übereinstimmung der Beobachtungen zu erhalten, habe ich die
Mühe nicht gescheut, den Betrag für jede einzelne Beobachtung zu berechnen. Dabei habe ich
einige kleine, sich leicht anbietende Rechenvorteile genutzt.

Die dritte Spalte enthält die Reduktion auf den mittleren Ort zu Jahresbeginn wegen
Aberration, Nutation und Präzession. Bei einigen Sternen kommt noch die Eigenbewegung hinzu.
Für die jährliche Eigenbewegung in Deklination wurde nämlich angenommen:
\begin{center}
\begin{tabular}{@{}lr@{\quad}c@{\quad}l@{}}
für&10&.\ .\ .\ .\ .\ .\ .&$-0''{,}42$\\
&25&.\ .\ .\ .\ .\ .\ .&$+0{,}33$\\
&37&.\ .\ .\ .\ .\ .\ .&$+0{,}38$.
\end{tabular}
\end{center}

\clearpage
\noindent\parbox[t]{\dimexpr\textwidth-10mm\relax}{\centering\scriptsize BESTIMMUNG DES BREITENUNTERSCHIEDS USW.\quad II.\quad DIE BEOBACHTUNGEN.}\makebox[10mm][r]{11}\par
\noindent Bei den beiden ersten Sternen steht die Eigenbewegung seit Langem fest. Bei Nr.\ 37
zeigt sie sich beim Vergleich mit der Bestimmung von \textsc{Piazzi} so, dass sie nicht zu
bezweifeln ist, sofern diese Bestimmung richtig ist\textsuperscript{*}). Der Berechnung
der Aberration, Nutation und Präzession liegen die wertvollen Tafeln von \textsc{Baily}
zugrunde. Nach ihnen wurde für jeden Stern mit Unterstützung der Herren von
\textsc{Nehus} und \textsc{Petersen} eine Ephemeride in Abständen von zehn Tagen berechnet
und darin unter Berücksichtigung der zweiten Differenzen interpoliert.

Die vierte Spalte enthält schließlich die Summe der ersten drei, also die wahre,
nur noch mit dem Kollimationsfehler behaftete Zenitdistanz für die mittlere Stellung
zu Beginn des Jahres 1827, wie sie sich aus jeder einzelnen Beobachtung ergibt.

\medskip\noindent\rule{.24\textwidth}{.4pt}\par
{\footnotesize
\noindent\textsuperscript{*}) Die Richtigkeit der auf 8 Beobachtungen beruhenden Bestimmung
dieses Sterns durch \textsc{Piazzi} wird durch die enge Übereinstimmung mit der Angabe in der
älteren Ausgabe seines Verzeichnisses von 1803 bestätigt, die auf 6 Beobachtungen beruhte.
Die genaue Größe der Eigenbewegung bleibt dennoch etwas ungewiss, weil das Jahr, das dem
Mittel der Beobachtungen entspricht, unbekannt ist. Es ist bemerkenswert, eine derart
beträchtliche Eigenbewegung bei einem Stern 7.\ Größe zu finden. Auch Nr.\ 11 scheint in
dieser Hinsicht die Aufmerksamkeit der Astronomen zu verdienen.\par}
\begin{center}\rule{.24\textwidth}{.4pt}\end{center}
